\documentclass[12pt]{article}
\usepackage[a4paper, margin=1in]{geometry}
\usepackage{amsmath, minted, enumitem, cmbright, hyperref}
\renewcommand{\thesection}{\Alph{section}}
\renewcommand{\thesubsection}{\thesection.\arabic{subsection}}
\renewcommand{\t}{\texttt}
\begin{document}
\section{Application Questions}
\subsection{Special Largest Gap}
As the numbers have at most $7$ digits, the digit sum is at most $7\times9=63$. Thus we can the keep the minimum/maximum seen so far for each possible digit sum in (fixed) lists and just scan the input once, giving $O(n)$ time complexity and $O(1)$ (extra) space complexity.

The idea is as follows:
\begin{itemize}
\item Maintain two lists of length $64$, $m$ and $M$, for the minimum and the maximum seen so far respectively.
\item For each number in the given array, compute its digit sum and update the corresponding element in $m$ and $M$.
\item Output the difference between the two numbers corresponding to the largest gap between $m$ and $M$, or $-1$ if not found.
\end{itemize}

A Python code is given below.
\begin{minted}{python}
class Solution:
    def specialLargestGap(self, arr: List[int]) -> int:
        def sod(n: int) -> int:
            s = 0
            while n:
                s += n % 10
                n //= 10
            return s
        m = [float('inf')] * 64
        M = [-1] * 64
        c = [0] * 64
        for i in arr:
            s = sod(i)
            m[s] = min(m[s], i)
            M[s] = max(M[s], i)
            c[s] += 1
        a = -1
        for s in range(64):
            if c[s] >= 2:
                a = max(a, M[s] - m[s])
        return a


tests = [[[63, 123, 27, 77, 7, 239, 45]],
         [[7, 16, 25, 34, 43, 52, 61, 70]],
         [[10, 2, 77, 14]]]
for t in tests:
    print(f"{t} -> {Solution().specialLargestGap(*t)}")
'''
[[63, 123, 27, 77, 7, 239, 45]] -> 162
[[7, 16, 25, 34, 43, 52, 61, 70]] -> 63
[[10, 2, 77, 14]] -> -1
'''
\end{minted}
\subsection{Linked List Splicing}
\textit{Note that the original problem is written in Java but I will do this in Python.}

The idea is as follows:
\begin{itemize}
\item Traverse \t{SLL1} to find the unique sublist \t{2 → 0 → 4 → 0}.
\item Denote \t{pre} as the node before $2$ and \t{aft} as the node after $0$.
\item Connect them to \t{SLL2} as \t{pre.next = head2} and \t{tail2.next = aft}.
\item Return \t{head1}.
\end{itemize}

A Python code is given below.
\begin{minted}{python}
# Definition for singly-linked list.
class ListNode:
    def __init__(self, val=0, next=None):
        self.val = val
        self.next = next


class Solution:
    def splice(self, head1: Optional[ListNode],
               tail1: Optional[ListNode],
               head2: Optional[ListNode],
               tail2: Optional[ListNode]) -> Optional[ListNode]:
        dum = ListNode(0, head1)
        pre = dum
        cur = head1
        while cur:
            if (cur.val == 2 and cur.next.val == 0 and cur.next.next.val == 4
                and cur.next.next.next.val == 0):
                pre.next = head2
                tail2.next = cur.next.next.next.next
                break
            pre = cur
            cur = cur.next
        return dum.next


def build(values):
    head = None
    for v in reversed(values):
        head = ListNode(v, head)
    return head


def to_tail(head):
    if not head:
        return None
    while head.next:
        head = head.next
    return head


def to_list(head):
    out = []
    while head:
        out.append(head.val)
        head = head.next
    return out


tests = [[[1, 1, 0, 1, 7, 2, 0, 4, 0, 7, 3, 2, 3, 0], [71, 72, 73, 74, 75]],
         [[2, 0, 4, 0, 9, 9], [8, 8, 8, 8]],
         [[5, 6, 2, 0, 4, 0], [1, 2, 3, 4]],
         [[9, 1, 2, 0, 4, 0, 3, 3, 3], [10, 20, 30, 40]]]

for SLL1, SLL2 in tests:
    head1 = build(SLL1)
    tail1 = to_tail(head1)
    head2 = build(SLL2)
    tail2 = to_tail(head2)
    _ = to_list(Solution().splice(head1, tail1, head2, tail2))
    print(f"{SLL1}, {SLL2} -> {_}")
'''
[1, 1, 0, 1, 7, 2, 0, 4, 0, 7, 3, 2, 3, 0], [71, 72, 73, 74, 75] ->
[1, 1, 0, 1, 7, 71, 72, 73, 74, 75, 7, 3, 2, 3, 0]
[2, 0, 4, 0, 9, 9], [8, 8, 8, 8] -> [8, 8, 8, 8, 9, 9]
[5, 6, 2, 0, 4, 0], [1, 2, 3, 4] -> [5, 6, 1, 2, 3, 4]
[9, 1, 2, 0, 4, 0, 3, 3, 3], [10, 20, 30, 40] -> [9, 1, 10, 20, 30, 40, 3, 3, 3]
'''
'''
\end{minted}
\subsection{Special Operations}
Notice that all elements are in $\left[7,7^7\right]$. In every operation we replace $x$ with $\max(\lfloor x/7\rfloor, 7)$, so $x$ becomes $7$ after at most $6$ operations. Thus all elements become $7$ after at most $6n$ operations and any remaining operations can only add exactly $7$ each.

The idea is to use a greedy algorithm. Indeed we can always pick the current maximum element. Otherwise if we pick any non-maximum element $y<x$, swapping to pick $x$ now and $y$ later can only increase the \t{total} as the values cannot go up.

Thus we can:
\begin{itemize}
\item Maintain a max-heap of the $n$ numbers.
\item Pop $x$, add to \t{total} and push $\max(\lfloor x/7\rfloor, 7)$ back. Repeat for $m$ times.
\item If $m>6n$, add another $7(m-6n)$ to \t{total}.
\item Return \t{total}.
\end{itemize}

Building the heap takes $O(n)$ and each operation takes $O(\log n)$. Thus the time complexity would be $O(n+\min(m,6n)\log n)=O(n\log n)$.

A Python code is given below.
\begin{minted}{python}
class Solution:
    def specialOperations(self, arr: List[int], m: int) -> int:
        from heapq import heapify, heappush, heappop
        n = len(arr)
        arr = list(map(lambda x: -x, arr))
        heapify(arr)
        total = 0
        for _ in range(min(m, 6 * n)):
            x = -heappop(arr)
            total += x
            heappush(arr, -max(x // 7, 7))
        return total + 7 * (m - 6 * n) if m > 6 * n else total


tests = [[[49, 49, 49], 4],
         [[441, 7, 8, 34], 4],
         [[441, 7, 8, 34], 77]]
for t in tests:
    print(f"{t} -> {Solution().specialOperations(*t)}")
'''
[[49, 49, 49], 4] -> 154
[[441, 7, 8, 34], 4] -> 547
[[441, 7, 8, 34], 77] -> 1059
'''
\end{minted}
\end{document}